\subsection{同型模}

设 A 是环，S 是一个单左 A-模。设 D 是 S 的自同态环的反环；它是一个体。赋以 A 与 D 的作用，S 是一个 (A,D)-双模。

命题 1. —— 设 M 是一个 A-模。下列性质等价：

\begin{enumerate}
\def\labelenumi{(\roman{enumi})}
\item
  存在一个集合 I，使 M 同构于 \(\mathrm{S}^{(\mathrm{I})}\) 。
\item
  模 M 是同构于 S 的子模的一个族的直和。
\item
  模 M 是同构于 S 的子模的一个族之和。
\item
  存在体 D 上的一个左向量空间 V，使 A-模 M 同构于 \(S \otimes_{D} V\) 。
\end{enumerate}

(i) 与 (ii) 的等价是直接的，而 (ii) 与 (iii) 的等价由 VIII，第 56 页的定理 1（应用于 N = 0 的情形）得出。D 上的每个左向量空间都同构于形如 \(\mathrm{D}_{s}^{(\mathrm{I})}\) 的向量空间，其中 I 是一个集合（II, §7, No.~1, 第 292 页，定理 1）。由于张量积与直和交换，(i) 等价于 (iv)。

定义 2. —— 若 A-模 M 具有命题 1 的诸等价性质，则称它是 S 型同型模。若存在一个单 A-模 T 使 M 是 T 型同型模，则称模 M 为同型模。

每个同型模都是半单的。

命题 2. —— 若一个模是 S 型同型子模之和，则它是 S 型同型模。S 型同型模的子模与商模都是 S 型同型的。

第一个断言由定义得出，第二个由 VIII，第 56 页的推论 3 得出。

注. —— 每个非零的 S 型同型模都有一个同构于 S 的商模与一个同构于 S 的子模；因此，若 M 与 M' 是非零的 S 型同型 A-模，则群 \(\operatorname{Hom}_{\mathrm{A}}(\mathrm{M},\mathrm{M}^{\prime})\) 不化为 0。

命题 3. —— a) 设 M 是一个 S 型同型 A-模。A-线性映射 \(\alpha_{\mathrm{M}}: \mathrm{S} \otimes_{\mathrm{D}} \mathrm{Hom}_{\mathrm{A}}(\mathrm{S}, \mathrm{M}) \to \mathrm{M}\)——它由 \(\alpha_{\mathrm{M}}(s \otimes f) = f(s)\)（VIII，第 59 页）刻画——是双射。

\begin{enumerate}
\def\labelenumi{\alph{enumi})}
\setcounter{enumi}{1}
\tightlist
\item
  设 V 是体 D 上的一个左向量空间。D-线性映射 \(\beta_{\mathrm{V}}: \mathrm{V} \to \mathrm{Hom}_{\mathrm{A}}(\mathrm{S}, \mathrm{S} \otimes_{\mathrm{D}} \mathrm{V})\)——它由 \(\beta_{\mathrm{V}}(v)(s) = s \otimes v\)（VIII，第 59 页）定义——是双射。
\end{enumerate}

把左 D-向量空间 \(\mathrm{Hom}_{\mathrm{A}}(\mathrm{S},\mathrm{M})\) 记作 \(\mathcal{H}(\mathrm{M})\)。按假设，A-模 M 是一族同构于 S 的子模的直和。A-模 S 是单生成的；为证映射 \(\alpha_{\mathrm{M}}\) 是双射，只需考虑 \(\mathrm{M} = \mathrm{S}\) 的情形（VIII，第 60 页，注 3）。而 \(\mathcal{H}(\mathrm{S})\) 就是 D-向量空间 \(\mathrm{D}_s\)，\(\alpha_{\mathrm{S}}\) 就是同构 \(\iota : \mathrm{S} \otimes_{\mathrm{D}} \mathrm{D}_s \to \mathrm{S}\)——它由 \(\iota(s \otimes d) = sd\) 定义。

同样，为证 b)，只需考虑 \(V = D_{s}\) 的情形。由于映射 \(\alpha_{S}\) 是双射，映射 \(\beta_{D_{s}} = \beta_{\mathcal{H}(S)}\) 也是双射（VIII，第 60 页，注 2）。

